\chapter{Introduction and approach}
\label{ch:sq1:intro}

This document answers the supplementary question on Power Usage Effectiveness reporting issued to Protocol Policy Lab following the hearing of 1 May 2026. The inquiry to which it contributes is itself a novel exercise, described by its proponents as a process to \enquote{establish a parliamentary inquiry into data centre developments in NSW, the first of its kind in Australia} \autocite{greensNSWDataCentreInquiry2026}. That novelty matters, because it means New South Wales is forming its evidentiary and regulatory baseline at the same moment that the sector is expanding fastest.

\section{Evidence base and method}
\label{sec:sq1:method}

The answer draws on two bodies of evidence. The first is the peer-reviewed and technical literature on data centre energy, water and grid integration. The second is the primary and grey literature specific to the Australian National Electricity Market and to New South Wales, including the Australian Energy Market Operator forecasts as relayed in the New South Wales Government's own consultation paper and analysis by industry and academic commentators. Quotations throughout are reproduced verbatim from the cited source. Measured quantities are distinguished from projections, and projections are attributed to their modelling scenario wherever the source specifies one. Where no New South Wales figure exists, this is stated rather than inferred from overseas data.

\section{Power Usage Effectiveness and the limits of a single metric}
\label{sec:sq1:pue-defn}

Power Usage Effectiveness is the ratio of a facility's total energy to the energy delivered to its information technology equipment. A value of 1.0 is the theoretical ideal in which no energy is spent on overheads, and lower values denote greater efficiency. The metric is mature and standardised, having been established such that \enquote{the assessment method of Power Usage Effectiveness has become an international standard for measuring energy efficiency and is widely used in the design and operation of data centers} \autocite{uanovAnalysisMethodsAssessing2022}. Cooling is the dominant overhead the ratio captures, since \enquote{the energy consumption of the cooling system for the machine room accounts for a significant part of the operating costs of the building} \autocite{uanovAnalysisMethodsAssessing2022}.

PUE is a necessary instrument but an insufficient one, and this submission is careful not to treat it as a complete measure. It is a ratio, so it is silent on absolute consumption, on the carbon intensity of the electricity drawn, and on water. The water dimension is consequential and under-measured, given that \enquote{Data centres consume water directly for cooling, in some cases 57\% sourced from potable water} \autocite{myttonDataCentreWater2021}, and that the academic literature records how, for data centres, \enquote{their water footprint (WF) has so far received little to no attention} \autocite{risticWaterFootprintData2015}. Water intensity is also enormously site-dependent, with one assessment finding \enquote{workload-level water use variations exceeding 10,000-fold} \autocite{leiWaterUseData2025}. A PUE figure can improve while absolute energy, emissions and water all rise, which is why the recent literature advocates \enquote{mandatory resource transparency, efficiency regulation, and lifecycle accountability} rather than a single efficiency ratio \autocite{imanHiddenCostsIntelligence2025}.

The insufficiency is compounded by rebound. Efficiency gains at the facility level are routinely offset by growth in demand, so that \enquote{the energy-intensive nature of digital infrastructure may conversely increase emissions}, with the gains from efficiency repeatedly overtaken by rising consumption \autocite{liuImpactGlobalDigital2025}. The United States experience illustrates the same point in reverse, where for two decades \enquote{comparable increases in electricity demand have been avoided through the adoption of key energy efficiency measures and a shift towards large cloud-based service providers} \autocite{shehabiDataCenterGrowth2018a}, a stabilisation that the artificial intelligence load is now ending. The practical implication for New South Wales is that an efficiency-conditioned regime should require a basket of measured indicators, being PUE alongside water usage effectiveness, absolute energy and renewable share, not PUE in isolation.

\section{The standard of inference adopted}
\label{sec:sq1:inference}

This submission's recommendation rests in part on inference rather than direct local measurement, and the standard applied is stated plainly so that the Committee can weigh it appropriately. The argument is that measurement to a common standard and disclosure are a precondition for any efficiency-conditioned or capacity-allocation policy, since a standard that is not measured cannot be enforced and a standard that is not disclosed cannot be scrutinised.
